\scanpage{48}
% archive sheet number: 47
\hfill\handwritten{6.1.1966}\\
Dear Coates,\\

Here a\handwritten{$\vert$}few more comments to my talk on the conjectures. The following
proposition shows that the conjecture $C_{\unsure{\ell|l|\chi}}(X)$ is independant of the choosen
polarization, and has also some extra interest, in showing the part playd
by the fact that $H^{i}(X)$ should be ``motive-theor\deleted{\unsure{i}}etically'' isomorphic to
its natural dual $H^{2n-i}(X)$ (as usual, I drop the twist for simplicity).

\uline{Proposition}\quad The condition $C_{\unsure{1|l}}(X)$ is equivalent also to \inserted{each of} the follwing
conditions:

a)\quad $D_{\ell}(X)$ ho\deleted{r}lds, and for every $i\inserted{\handwritten{<}n}$, there exists an isomorphism
$H^{2n-i}(X)\handwritten{\to}H^{i}(X)$ which is algebraic (i.e. induced by an algebraic correspondance
class; we do not make any assertion on what it induces in degrees
different from $2n-i$).

b)\quad For every endomorphism $H^{i}(X)\handwritten{\unsure{-|\to}}H^{i}(X)$ which is algebraic, the
coefficient\unsure{s} of the characteristic polynomial are rational, and for every
$i\inserted{\handwritten{<}n}$, there exists an isomorphism $H^{2n-i}(X)\handwritten{-}H^{i}(X)$ which is algebraic.

Proof. I sketched already how $D_{\ell}(X)$ implies the fact that for an
algebraic endomorphism of $H^{i}(X)$, the coefficients of the characteristic
polynomial are rational numbers, Therefore we know that a) implies b),
and of course $C_{\ell}(X)$ implies a). It remains to prove that b) implies
$C_{\ell}(X)$. Let $u:H^{2n-i}(X)\handwritten{-}H^{i}(X)$ be the given isomorphism which is algebraic,
and $v:H^{i}(X)\handwritten{-}H^{\deleted{\unsure{a}}2n-i}(X)$ \deleted{an}\handwritten{the} algebraic isomorphism in the opposite
direction, induced by $L_{X}^{n-i}$. Then $uv=w$ is an automorphism of $H^{i}(X)$ which
is algebraic, a\deleted{\unsure{ad}}nd the Hamilton-Cayley formula $\unsure{u|w}^{\unsure{h|b}}-\handwritten{\sigma}_{1}(w)u^{h-1}+\mathinner{...}+\deleted{(-1)}$
$(-1)^{b}\handwritten{\sigma}_{b}(w)=0$ (where the $\handwritten{\sigma}_{i}(w)$ are the coefficients of the car. pol.
of $w$) whos that $w^{-1}$ is a linear combination of the $w^{i}$, with coe\deleted{\unsure{e}}fficient
of the type $\pm\handwritten{\sigma}_{i}(w)/\handwritten{\sigma}_{b}(w)$ (NB $b=\operatorname{rank}H^{i}$). The assumption implies
